% =============================================================================
% Partie B — Conception de l'application
% =============================================================================
\section{Partie B — Conception de l'application}

\subsection{Cas d'utilisation}

Trois acteurs, onze cas. La lecture utile n'est pas la liste mais les deux
relations entre cas, qui portent chacune une décision.

\illus{05-cas-utilisation}{0.88}{Cas d'utilisation : rédacteur, visiteur,
administrateur.}{fig:cas-utilisation}

La première relation paraît anodine et ne l'est pas : « éditer en collaboration »
\emph{étend} « consulter son crédit », et non l'inverse. Le crédit n'est jamais
interrogé avant d'écrire. Le retour du serveur — acquittement ou refus — porte
déjà le crédit restant, et c'est le refus qui met la jauge à jour. Une relation
d'inclusion aurait décrit un aller-retour qui n'existe pas.

La seconde est une absence qu'il faut assumer : il n'y a \emph{aucune} relation
entre « réserver un exemplaire » et « s'identifier ». Le visiteur ne s'identifie
jamais. Ce n'est pas un oubli dans la conception, c'est la décision qui rend le
parcours visiteur absorbable à 15\,000 clics : pas de session à créer, pas de
jeton à vérifier, une seule écriture conditionnelle.

\subsection{Diagramme de paquets}

L'organisation du code suit une règle unique : les routes dépendent du domaine,
le domaine dépend des accès aux services, jamais l'inverse. Un seul paquet parle
aux services AWS. Le domaine reçoit le stockage, la file et le courriel par
injection, ce qui permet aux tests de n'appeler aucun service.

\illus{06-packages}{0.80}{Organisation des paquets du backend.}{fig:packages}

Cette organisation n'est pas décorative. C'est elle qui a permis d'écrire
43 tests sans aucun appel réseau, et c'est elle qui a rendu possible le mode
local complet, où l'application tourne avec un stockage et une base en mémoire,
sans compte AWS.

\subsection{Diagramme de classes et modèle de données}

Le modèle tient en dix tables, dont trois portent une règle plutôt qu'une donnée.

\illus{07-classes}{0.86}{Modèle du domaine : dix entités, trois portent une
règle de concurrence ou de quota.}{fig:classes}

Trois points de conception méritent l'attention, parce qu'ils sont faciles à
modéliser faux.

Le premier concerne les \textbf{versions}. Le numéro de version \emph{est} le
numéro de séquence du document au moment de l'enregistrement. Ce n'est pas deux
champs mais un seul, utilisé deux fois. La conséquence pratique est agréable : la
numérotation ne peut pas avoir de trous, puisqu'elle suit une séquence déjà
garantie contiguë par le journal.

Le deuxième concerne l'\textbf{idempotence}. La table des marqueurs porte deux
clés par mise à jour : l'identifiant d'opération fourni par le client, et
l'empreinte du contenu. La seconde existe parce qu'un client qui se reconnecte
peut renvoyer la même modification sous un autre identifiant ; sans elle, un
renvoi débiterait le crédit deux fois.

Le troisième concerne le \textbf{cache de PDF}. Il porte un compteur de
tentatives et une date de bail. Sans eux, deux visiteurs arrivant sur un PDF
absent publieraient deux travaux de génération, et un PDF qui échoue serait
redemandé indéfiniment.

\begin{table}[H]
\centering
\small
\begin{tabular}{@{}p{3.6cm}p{5.2cm}p{6.6cm}@{}}
\toprule
\textbf{Table} & \textbf{Clé} & \textbf{Rôle de conception} \\
\midrule
\code{Documents} & \code{docId} & séquence courante, dernier instantané, longueur du
  texte \\
\code{Members} & \code{docId} + \code{email} & appartenance ; index secondaire par
  adresse pour lister les documents d'un rédacteur \\
\code{Updates} & \code{docId} + \code{seq} & journal : une entrée par lot, avec
  l'update Yjs fusionné \\
\code{Versions} & \code{docId} + \code{n} & instantané consultable ; \code{n} est
  le numéro de séquence \\
\code{Links} & \code{token} & quota : \code{remaining} et statut \\
\code{Tickets} & \code{token} + \code{visitorId} & exemplaire réservé et numéro
  figé \\
\code{PdfCache} & \code{docId\#seq} & état de génération, verrou, tentatives \\
\code{Credits} & \code{email\#jour} & compteur consommé et plafond \\
\code{Abus} & \code{email\#jour} & refus du jour et terme du blocage \\
\code{OpIds} & \code{id} & idempotence : identifiant d'opération et empreinte \\
\bottomrule
\end{tabular}
\caption{Modèle de données : dix tables, une clé et un rôle par table.}
\label{tab:tables}
\end{table}

\subsection{Diagrammes d'états-transitions}

Trois objets ont un cycle de vie qui porte une règle : le lien, le visiteur, et
le crédit.

\subsubsection{États du lien}

Un lien naît actif avec son quota, le consomme, et devient épuisé. La transition
est \emph{monotone} : un lien ne redevient jamais actif. Point important pour la
lecture du diagramme : l'état épuisé \textbf{n'est pas terminal}. Le lien reste
consultable indéfiniment et répond « épuisé », ce qui est précisément ce que
demande C6 — plus aucun travail inutile.

\illus{08a-etats-lien}{0.52}{États d'un lien de téléchargement.}{fig:etats-lien}

\subsubsection{États du visiteur}

Le visiteur suit un chemin qui dépend de deux choses indépendantes : le quota du
lien et l'état du fichier. C'est ce qui produit trois réponses distinctes là où
on attendrait une seule : prêt, patienter, épuisé.

\illus{08b-etats-visiteur}{0.62}{États d'un visiteur sur un lien.}{fig:etats-visiteur}

Un point mérite d'être souligné : une fois l'exemplaire réservé, il est
\emph{acquis}. L'URL présignée expire au bout de soixante secondes, mais le
visiteur peut en redemander une indéfiniment. Rien de ce qu'il fait ensuite ne
consomme d'exemplaire supplémentaire.

\subsubsection{États du crédit}

Le crédit a trois états, et le troisième est celui qu'on oublie : un utilisateur
peut être \emph{bloqué alors qu'il lui reste du crédit}. Le blocage est posé par
quatre refus dans la journée ; il ne s'efface pas parce que le compteur a été
remis à zéro.

\illus{08c-etats-credit}{0.62}{États du crédit d'écriture, par jour et par
adresse.}{fig:etats-credit}

La remise à zéro mérite une explication, parce qu'elle a été source de confusion.
Ce n'est pas un événement, c'est un changement de clé : le compteur est indexé
par adresse et par jour. Deux journées différentes sont deux lignes différentes,
sans qu'aucune écriture ne soit nécessaire.

\hypothese{Le « jour » est celui de Douala, pas le jour UTC. C'est un écart avec
l'hypothèse initiale du plan. Un utilisateur lisait « recharge à 01:00 » et ne
comprenait pas : minuit UTC tombe à une heure du matin sur place. Le décalage
est constant — le Cameroun n'applique pas d'heure d'été — ce qui rend
l'arithmétique exacte toute l'année. Le terme du blocage est par ailleurs
\emph{recalculé} à chaque lecture et non relu en base, ce qui a corrigé les
lignes écrites sous l'ancienne règle sans aucune migration.}

\subsection{Diagrammes de séquence système}

Trois parcours structurent l'application : réserver puis télécharger, éditer,
restaurer.

\subsubsection{Réserver et télécharger}

\illus{09-ssd-telechargement}{0.80}{Séquence système : réserver puis télécharger
un exemplaire.}{fig:ssd-telechargement}

Deux précisions que le diagramme ne peut pas montrer. Consulter l'état du lien
ne consomme rien : c'est une exigence de conception, sans quoi un aperçu de lien
dans une messagerie ou un robot d'indexation épuiserait le quota. Et
l'interrogation d'attente porte sur le \emph{fichier}, pas sur l'état du lien :
l'état du lien ne dit rien de l'avancement de la génération.

\subsubsection{Éditer un document}

\illus{10-ssd-ecriture}{0.72}{Séquence système : éditer un document en
collaboration.}{fig:ssd-ecriture}

Le comportement en cas de refus mérite d'être explicite, parce qu'il est
contre-intuitif. À la réception d'un refus de crédit, le client vide sa file
locale et reconstruit le document depuis le serveur. Ce n'est pas une
resynchronisation : le refus n'a rien modifié. Le client reconstruit parce que
les frappes déjà saisies ne seront \emph{jamais} enregistrées, et qu'il faut le
montrer sans ambiguïté plutôt que laisser un état intermédiaire trompeur.

\subsubsection{Restaurer une version}

\illus{11-ssd-restauration}{0.78}{Séquence système : restaurer une version.}{fig:ssd-restauration}

La restauration est une écriture, pas un retour en arrière. Trois propriétés en
découlent, et toutes les trois sont voulues : la compatibilité avec le CRDT, donc
la convergence des autres clients ; la conservation de l'historique, puisque le
texte écrasé est d'abord sauvegardé comme version ; la gratuité, puisqu'aucun
caractère n'est facturé et que l'opération franchit un blocage. Cette dernière
propriété est un choix : sans elle, un utilisateur ayant épuisé son crédit serait
enfermé sans recours.

\subsection{Conception par contrainte}

Cette section reprend les six contraintes une par une, et indique pour chacune
\emph{où} elle se tient dans la conception, et \emph{comment} elle est vérifiée.
C'est la lecture la plus utile du dossier : elle relie chaque exigence à un
mécanisme et à une preuve.

\subsubsection{C1 — exactement $N$ téléchargements}

\textbf{Où.} Dans une seule transaction conditionnelle sur la table des liens et
la table des tickets. Le décrément n'a lieu que si le compteur est positif et le
lien actif ; le ticket n'est créé que si ce visiteur n'en a pas.

\textbf{Ce qui rend la faute possible, et comment elle est évitée.} Quatre
pièges ont été identifiés :

\begin{enumerate}[leftmargin=1.4em,itemsep=2pt]
  \item chercher le ticket existant \emph{après} le court-circuit « épuisé » :
        un visiteur déjà servi recevrait alors un refus. L'ordre est inversé
        exprès ;
  \item compter l'exemplaire au téléchargement du fichier et non à la
        réservation : un visiteur qui recharge la page consommerait deux fois.
        L'exemplaire est consommé au claim ;
  \item poser la transition « épuisé » dans la même écriture que le décrément :
        cinquante claims simultanés tenteraient alors cinquante fois la même
        transition. Elle est séparée et conditionnelle ;
  \item croire que la garde en mémoire décide du quota. Elle ne fait
        qu'\emph{alléger} la base sous le pic ; un visiteur qu'elle écarte reçoit
        « patienter », jamais « épuisé ».
\end{enumerate}

\textbf{Preuve.} Onze tests unitaires, dont un test de concurrence : cinquante
claims simultanés pour vingt exemplaires donnent exactement vingt réservations.
En charge, 15\,000 visiteurs sur un lien de 500 exemplaires ont produit
\emph{exactement} 500 réservations, sans une seule erreur serveur.

\subsubsection{C2 — aucune modification acceptée perdue}

\textbf{Où.} Dans l'écrivain unique par document et dans l'ordre
acquittement-après-écriture. Un lot est écrit en une transaction atomique, puis
diffusé en un seul message. Un acquittement ne part qu'après validation de la
transaction.

\textbf{Ce qui rend la faute possible, et comment elle est évitée.} Trois
risques :

\begin{enumerate}[leftmargin=1.4em,itemsep=2pt]
  \item acquitter avant d'écrire : une panne entre les deux perd la modification
        sans que personne ne le sache. L'acquittement est émis après ;
  \item renvoyer une modification sous un nouvel identifiant après une coupure :
        elle serait écrite deux fois. Le client renvoie le \emph{même}
        identifiant, et le serveur déduplique aussi par empreinte du contenu ;
  \item mesurer la convergence avec un outil qui ne peut pas la mesurer. Un
        générateur de charge compte des acquittements, il ne fait pas tourner un
        client CRDT.
\end{enumerate}

\textbf{Preuve.} Deux preuves distinctes, et c'est le point important. Le
générateur de charge vérifie que chaque modification émise reçoit son
acquittement. Un second outil, écrit en Node avec un vrai client CRDT,
reconstruit le document depuis le journal, vérifie que les numéros de séquence
sont contigus, et compare deux reconstructions indépendantes par empreinte. Ce
second outil a été validé dans les deux sens : il conclut à la convergence sur un
journal sain, et à l'échec dès qu'un numéro manque. Un détecteur qu'on n'a jamais
vu détecter ne prouve rien.

\subsubsection{C3 — document tel qu'il est à $T_1$}

\textbf{Où.} Au moment du claim : le dernier numéro de séquence réellement écrit
est figé dans le ticket. Le PDF est ensuite construit depuis le journal jusqu'à
ce numéro.

\textbf{Ce qui rend la faute possible.} Le PDF pourrait être généré à la demande,
au moment du téléchargement : deux visiteurs d'une même rafale recevraient alors
des documents différents, ce qui viole l'exigence de complétude et de cohérence.
Il pourrait aussi être régénéré à chaque visite, ce qui multiplierait le travail
par le nombre de visiteurs au lieu du nombre de versions.

\textbf{Preuve.} Le contenu du PDF est déterminé par le numéro figé, et la clé de
cache inclut ce numéro : une version, un fichier. Un verrou avec bail empêche
deux générations simultanées pour la même version.

\subsubsection{C4 — jamais plus que son crédit}

\textbf{Où.} Dans un compteur par adresse et par jour, débité dans la même
transaction que l'écriture au journal. Le comptage est effectué sur le binaire
Yjs, côté serveur : le client ne déclare pas ce qu'il a écrit.

\textbf{Ce qui rend la faute possible, et comment elle est évitée.}

\begin{enumerate}[leftmargin=1.4em,itemsep=2pt]
  \item compter côté client : un client modifié déclarerait zéro. Le comptage
        est serveur, en décodant la mise à jour ;
  \item ne pas facturer les suppressions : écrire puis effacer serait gratuit, et
        le plafond contournable en trois gestes. Les suppressions sont facturées ;
  \item compter par connexion plutôt que par adresse : cinq onglets donneraient
        cinq quotas. La clé est l'adresse ;
  \item inscrire le plafond comme condition de la transaction : DynamoDB refuse
        l'arithmétique entre valeurs. C'est une lecture préalable.
\end{enumerate}

\textbf{Preuve.} Seize tests unitaires, dont l'acceptation exacte jusqu'au
plafond et pas un caractère de plus, le partage du compteur entre plusieurs
connexions, la déduplication par identifiant et par empreinte, et le blocage au
quatrième refus.

\subsubsection{C5 — 15\,000 clics et 300 rédacteurs}

\textbf{Où.} À quatre endroits complémentaires : le micro-cache du serveur
frontal, la garde en mémoire sur les réservations, la diffusion par lots, et le
dimensionnement de l'instance.

\textbf{Ce qui rend la faute possible.} Chaque mécanisme a un mode de défaillance
propre, et c'est leur combinaison qui tient :

\begin{itemize}[leftmargin=1.4em,itemsep=2pt]
  \item sans micro-cache, 15\,000 lectures d'état frappent la base et la font
        échouer en limitation de débit ;
  \item sans garde en mémoire, autant d'écritures conditionnelles simultanées
        produisent le même effet ;
  \item sans diffusion par lots, un message par frappe multiplie la charge
        réseau par le nombre de rédacteurs ;
  \item sans dimensionnement mesuré, une instance trop petite s'effondre
        précisément pendant le pic qu'on mesure.
\end{itemize}

\textbf{Preuve.} Le scénario mixte fait tourner visiteurs, rédacteurs et sondes
d'état en même temps. Les mesures relevées sont la latence de réservation, la
latence d'état, la latence d'acquittement, et le nombre d'erreurs serveur —
celui-ci doit rester à zéro, y compris pour le parcours visiteur.

Un piège de mesure mérite d'être signalé : une réponse « épuisé » est un code 410,
que les outils de charge comptent comme une requête échouée. Un seuil sur ce taux
d'échec est donc inatteignable par construction — 84\,\% mesurés sans la moindre
panne. Le critère utile est un compteur dédié aux erreurs serveur.

\subsubsection{C6 — dernier exemplaire : prévenir vite, arrêter le travail}

\textbf{Où.} Dans la monotonie de l'état épuisé et dans son cache. Une fois le
lien épuisé, la réponse est immédiate, ne touche pas la base, et le serveur
frontal la conserve longuement. Le créateur du lien reçoit un courriel unique.

\textbf{Ce qui rend la faute possible.} Deux choses : recalculer l'état à chaque
visite, ce qui revient à travailler pour ne rien dire ; et envoyer le courriel
plusieurs fois, une par visiteur qui constate l'épuisement. Le courriel est donc
protégé par un verrou conditionnel : seul le claim qui pose la transition
déclenche l'envoi.

\textbf{Preuve.} L'état épuisé est posé une seule fois, et la notification est
idempotente : neuf tests unitaires couvrent l'envoi au seul dernier exemplaire,
l'absence de second envoi, la panne du service de courriel sans impact pour le
visiteur, et le relâchement du verrou en cas d'échec — mieux vaut un doublon rare
qu'un créateur jamais prévenu.

\begin{table}[H]
\centering
\small
\begin{tabular}{@{}p{0.7cm}p{5.6cm}p{8.7cm}@{}}
\toprule
\textbf{\#} & \textbf{Mécanisme} & \textbf{Preuve exécutable} \\
\midrule
C1 & transaction conditionnelle, ordre des vérifications, transition séparée
   & 11 tests dont concurrence 50/20 ; en charge, 500 réservations exactes \\
\addlinespace
C2 & écrivain unique, acquittement après écriture, déduplication par empreinte
   & acquittements sous charge ; contiguïté du journal et double reconstruction \\
\addlinespace
C3 & numéro figé au claim, une génération par version, verrou avec bail
   & clé de cache par version ; un seul travail par version \\
\addlinespace
C4 & comptage serveur sur le binaire, suppressions facturées, clé par adresse
   & 16 tests : plafond exact, multi-connexions, blocage au 4\up{e} refus \\
\addlinespace
C5 & micro-cache, garde mémoire, diffusion par lots, dimensionnement mesuré
   & scénario mixte : visiteurs, rédacteurs et sondes simultanés \\
\addlinespace
C6 & état épuisé monotone et mis en cache, notification idempotente
   & 9 tests : envoi unique, panne sans impact, verrou relâché \\
\bottomrule
\end{tabular}
\caption{Synthèse : une contrainte, un mécanisme, une preuve.}
\label{tab:synthese-contraintes}
\end{table}
